(1) 颈管空气以相同速度运动，有等效质量 $m_e=\rho S\ell_e$。空气绝热体积模量为 $\rho c^2$，颈部位移引起腔体体积变化 $\Delta V=Sx$，压强变化 $\Delta p=-\rho c^2Sx/V$，颈管受回复力 $S\Delta p=-\rho c^2S^2x/V$，所以弹簧系数 $K_e=\rho c^2S^2/V$。

(2) 自由振动 $m_e\ddot x+K_ex=0$，因此 $\boxed{\omega_H=c\sqrt{S/(V\ell_e)}}$、$f_H=\omega_H/(2\pi)$。频率中的密度因子在空气颈惯性与腔体压缩刚度中相消，但声速一般会随温度改变。

(3) 受驱方程为 $m_e\ddot x+b\dot x+K_ex=Sp_a\cos\omega t$，稳态位移幅值 $A=Sp_a/\sqrt{(K_e-m_e\omega^2)^2+b^2\omega^2}$。若 $b\ll2m_e\omega_H$，位移共振附近的品质因数为 $\boxed{Q\simeq m_e\omega_H/b}$，带宽约 $\omega_H/Q$。

(4) 总机械能 $E=m_e\dot x^2/2+K_ex^2/2$，阻尼平均功率 $\bar P=b\omega^2A^2/2$，等于外部驱动的周期平均声功率。$S$ 增大时 $\omega_H\propto\sqrt S$，有效颈长变长或腔体积增大时频率下降。开口声流延伸到颈外，会以端部修正增加 $\ell_e$，小孔近似要求波长远大于腔和颈尺寸。
